\begin{table}[H]
\caption{Feature-based detector families as geometry readouts. The entries
summarize the framework only; they do not report experimental results.}
\label{tab:readout-map}
\begin{center}
\footnotesize
\setlength{\tabcolsep}{3pt}
\renewcommand{\arraystretch}{1.15}
\begin{tabular}{@{}>{\raggedright\arraybackslash}p{0.17\textwidth}>{\raggedright\arraybackslash}p{0.25\textwidth}>{\raggedright\arraybackslash}p{0.27\textwidth}>{\raggedright\arraybackslash}p{0.20\textwidth}@{}}
\toprule
Detector family & Score form & Geometry channel & Diagnostic control \\
\midrule
Mahalanobis
& \(-\min_c (f-\mu_c)^\top \Sigma_W^{-1}(f-\mu_c)\)
& class centers, within-class covariance, precision spectrum
& post-hoc L2 normalization; covariance shrinkage \\
Gaussian density
& \(\log\sum_c \pi_c \mathcal{N}(f\mid\mu_c,\widehat{\Sigma}_c)\)
& Gaussian density, covariance volume, anisotropy, conditioning
& tied, diagonal, and shrinkage covariance \\
kNN
& \(-d_k(f,\mathcal{B})\)
& local feature-bank spacing, radial scale, angular neighborhoods
& raw kNN vs. L2-kNN \\
CTM angular readout
& \(\max_c \langle \widehat f,a_c\rangle\)
& prototype or classifier directions, angular class separation
& normalized features; prototype vs. classifier direction \\
NECO collapse/residual readout
& \(\max_c\langle\widehat f,\widehat\mu_c\rangle\), \(-\|(I-P_M)f\|_2^2\)
& collapse-like prototype geometry, residual energy outside class-mean subspace
& prototype vs. residual readout; subspace dimension \\
\bottomrule
\end{tabular}
\end{center}
\end{table}
